% ============================================
% DOC-24
% Extraído de main.pandoc.tex
% ============================================
\fragtitle{doc}{D}{DOC}{24}{Parecer Técnico de Conservação}
\markboth{DOC-24 --- Parecer Técnico de Conservação}{}

\noindent \textit{Documento extraído do arquivo pericial.}

\vspace{0.5em}\noindent\rule{\textwidth}{0.5pt}\vspace{0.5em}

\noindent {\footnotesize\texttt{ARQUIVO P\'UBLICO $\cdot$ PROVA DOCUMENTAL $\cdot$ PROCESSO 1.260}}\\[-0.2em]

\noindent \textbf{Arquivo Público Mineiro --- Laboratório de Conservação e Restauro}\\
\noindent \textbf{Objeto:} caderno encadernado, registro 1.260.\\
\noindent \textbf{Técnica responsável:} Fernanda Lopes Bastos, CRT-04/1128.\\
\noindent \textbf{Solicitação:} avaliação de integridade física e recomendação de acondicionamento.\\
\noindent\textbf{Data:} 04 de agosto de 2026.

\noindent \textbf{1. Descrição geral}

\noindent Caderno em papel de gramatura variável (entre 80 e 120\,g/m²), encadernação artesanal em couro curtido, provavelmente produção local ou artesanal. Estado geral de conservação é inconsistente com a idade declarada (datas entre 1917 e 1979). O papel não apresenta os sinais esperados de oxidação para um objeto do período. A acidez do suporte é neutra --- incomum para papel de produção não industrial do início do século XX.

\noindent \textbf{2. Achado principal}

\noindent Durante a análise de fólios, foi identificada uma discrepância entre a espessura do caderno e a quantidade de páginas registradas no inventário de 1979. O inventário original --- elaborado pelo Dr. Antônio Tavares no momento da exumação --- lista aproximadamente 190 páginas preenchidas. A contagem atual indica 208 fólios com texto.

\noindent Não há evidência de que páginas tenham sido adicionadas fisicamente. A encadernação está íntegra, sem sinais de desmonte ou reinserção. Costura original, cola original. E, no entanto, há 18 páginas a mais.

\noindent \textbf{3. Análise das páginas excedentes}

\noindent As 18 páginas adicionais estão distribuídas ao longo de todo o caderno, intercaladas entre os fólios originais. Em cada caso, a página nova --- ou aparentemente nova --- insere-se entre duas páginas cujo texto é contínuo e cuja caligrafia corresponde ao período imediatamente anterior. A página adicional parece ter ``crescido'' entre as duas, como se o caderno tivesse gerado papel novo para comportar texto que não existia na escrita original.

\noindent A tinta das páginas adicionais é da mesma composição química da tinta original, mas o teor de ferro na formulação é ligeiramente superior --- indicando oxidação mais recente. O teste de datação relativa posiciona a escrita dessas páginas entre 2020 e 2026.

\noindent As páginas adicionais foram escritas no interior de um caderno lacrado em cofre desde 1979.

\noindent \textbf{4. Conteúdo das páginas adicionais}

\noindent \textit{A pedido da diretoria, o conteúdo das 18 páginas não foi transcrito. Informo apenas que a letra não corresponde à caligrafia do paciente 1.260. A caligrafia das páginas adicionais assemelha-se --- preliminarmente, sem análise grafotécnica conclusiva --- à do Dr. Heitor Oliveira em seus últimos registros.}
\noindent \textbf{5. Recomendação técnica}

\noindent Acondicionamento em envelope selado com atmosfera de argônio. Monitoramento semanal de peso e espessura. Proibição de manuseio individual. A técnica responsável solicita --- por meio deste parecer --- não ser convocada para nova análise do objeto.

\noindent \textit{Fernanda Lopes Bastos}\\
\noindent \textit{Técnica em Conservação e Restauro}\\
\noindent \textit{CRT-04/1128}

\vspace{0.5em}
\noindent\textit{\textbf{↪ Rotas:} \rota{DOC-01} \rotasep \rota{LUC-02} \rotasep \rota{DOC-25}}
